%%% ---------------------------------------------------------------------------
%%% preamble.tex —— 你自己的宏包与自定义命令
%%%
%%% 主题里已经载入了：
%%%   fontspec / unicode-math / acadmath / amsmath / mathtools
%%%   graphicx / booktabs / array / translator / appendixnumberbeamer
%%% 不要重复载入。特别注意：**不要载入 enumitem**，它会破坏 \item<2-> 覆盖语法。
%%% ---------------------------------------------------------------------------

%%% ---- 参考文献 ----
%%% 样式与 article 模板保持一致；投英文会议时换成 style=ieee
\usepackage[
  backend     = biber,
  style       = gb7714-2015,   % gb7714-2015 | gb7714-2015ay | ieee | numeric-comp
  maxbibnames = 3,
  minbibnames = 3,
  giveninits  = true,
  doi         = false,
  isbn        = false,
  url         = false,
]{biblatex}

%%% ---- 节页：每进入一个新 \section 自动插一页过渡 ----
%%% 不想要就把整段注释掉。
\AtBeginSection[]{%
  \begin{frame}[plain, noframenumbering]
    \sectionpage
  \end{frame}%
}

%%% ---- 按需追加的宏包 ----
% \usepackage{tikz}
% \usepackage{pgfplots}\pgfplotsset{compat=1.18}
% \usepackage{multicol}

%%% ---- 本报告专用的记号（与论文保持一致）----
\newcommand{\ourmethod}{MSFF-Net}
\newcommand{\dataset}{DOTA-v2.0}
\newcommand{\featmap}{\mat{F}}
\newcommand{\loss}{\mathcal{L}}
\newcommand{\params}{\vect{\theta}}

%%% ---- 占位图 ----
%%% 只为让模板在没有图片时也能编译；真正插图时改用
%%% \includegraphics[width=0.8\linewidth]{architecture}
\newcommand{\placeholder}[3]{%
  \begingroup
    \fboxsep=0pt \fboxrule=0.4pt
    \fbox{\parbox[c][#2][c]{\dimexpr#1-2\fboxrule\relax}{\centering\footnotesize\ttfamily\color{AcadMuted}#3}}%
  \endgroup
}

%%% ---- 演讲备注 ----
%%% 在帧里写 \note{...} 记备注。
%%%   make        普通编译，备注不出现
%%%   make notes  生成双屏 PDF（左边幻灯片、右边备注），投影时用 pdfpc 播放
%%% 开关由 make 传进来的 \AcadShowNotes 控制，这里不用手改。
\ifdefined\AcadShowNotes
  \usepackage{pgfpages}
  \setbeameroption{show notes on second screen = right}
\else
  \setbeameroption{hide notes}
\fi
